\section{Security Analysis}\label{sec:security}

We analyze threats as attack paths across trust boundaries, not as proof that a
named control defeats every adversary.  The relevant assets are authorization
state, source and bytecode, policy results, provider selection, credentials and
secret material, trace/report integrity, and jurisdictional metadata.
Adversaries may supply malformed programs, attempt unauthorized operations,
tamper with stored artifacts, exploit a provider adapter, or induce a consumer
to accept a favorable summary while ignoring detailed failures.

\begin{figure}[t]
\centering
\includegraphics[width=\linewidth]{threat-mitigation.pdf}
\caption{Threat--mitigation mapping for the evaluated architecture. Mapping
means a control is present or declared; it is not a measured attack-blocking
rate.}
\label{fig:threat-mitigation}
\end{figure}

\begin{table}[t]
\centering
\caption{Principal threats and bounded mitigation claims.}
\label{tab:threats}
\begin{tabular}{@{}lll@{}}
\toprule
Threat & Mitigation & Claim status \\
\midrule
Unauthorized provider execution & Executable-provider allowlist & Implemented \\
Network side effects & Offline and network-denied profiles & Implemented \\
Secret disclosure & Secret and key-material denial & Implemented \\
Evidence tampering & Digest-linked artifact bundle & Implemented \\
\bottomrule
\end{tabular}

\end{table}

\Cref{fig:threat-mitigation} visualizes which boundaries a threat crosses, and
\cref{tab:threats} states the narrower mitigation claim supported at each
boundary; neither reports measured adversarial efficacy.

\paragraph{Unauthorized provider execution.}
The split between declarative and executable registries prevents a provider
contract from becoming callable merely by being named.  Capability checks,
approval requirements, and a simulated-only allowlist reduce the observed
attack surface.  A malicious executable adapter or VM defect remains in scope
for future testing.

\paragraph{Network and secret side effects.}
The demonstrated runtime profile denies network access, external execution,
plaintext secrets, and key material.  Events make these decisions visible.
Because this study does not instrument the host operating system, the evidence
supports VM-level denial behavior, not production-grade isolation.

\paragraph{Artifact tampering.}
Semantic digests detect changes relative to an EvidenceBundle and avoid
formatting-dependent hashes.  An attacker able to replace both artifacts and
the unsigned bundle can create a self-consistent replacement.  Signatures,
trusted timestamps, independent publication, and key governance are outside
the implementation evaluated here.

\paragraph{Policy confusion.}
The observed status contradiction is an integrity-of-interpretation threat.
A dashboard reading only \code{security_checks} could present an unsafe
conclusion even when the detailed object requires review.  Consumers should
apply a precedence rule: any explicit failure, unknown rule, or required review
prevents an aggregate pass.  Future schemas should compute one canonical
tri-state decision from typed component results.

\paragraph{Prompt injection and content attacks.}
Prompt injection is a recognized application risk \citep{owasp_llm_top10}.
\Argorix can constrain resulting operations through capabilities and policy,
but this corpus contains no prompt text and no controlled attacks.  We make no
claim of universal injection prevention and report no success rate.

\paragraph{Identity and trust spoofing.}
Local semantic checks prevent dangling identity and credential references.
They do not establish control of a DID, truth of a credential, or execution of
a cryptographic handshake.  A future resolver and verifier would introduce
new trust anchors, replay concerns, revocation, privacy, and network-failure
paths \citep{w3c_did_core_2022,w3c_vc_data_model_2025}.

The strongest supported security statement is therefore conditional: within
the observed runtime path, requests are mediated by compiled policy and an
executable-provider allowlist, denials are recorded, and complete artifacts are
offline-consistent.  Security outside that path remains unmeasured.
